\TieudeT{PHONG TỎA VẬT LÍ 12 - VẬT LÍ NHIỆT}
\subsubsection*{PHẦN I. Thí sinh trả lời câu hỏi từ câu 1 đến câu 18. Mỗi câu hỏi thí sinh chỉ chọn một phương án}
\setcounter{ex}{0}
\begin{ex}
Nhiệt dung riêng của rượu là $2500\ J/kg.K$. Điều đó có nghĩa là
\choice
{nhiệt lượng có trong $1\ kg$ chất ấy ở nhiệt độ bình thường}
{để nâng $1\ kg$ rượu lên nhiệt độ bay hơi ta phải cung cấp cho nó một nhiệt lượng là $2500\ J$}
{$1\ kg$ rượu bị đông đặc thì giải phóng nhiệt lượng là $2500\ J$}
{\True để nâng $1\ kg$ rượu tăng lên $1$ độ ta cần cung cấp cho nó nhiệt lượng là $2500\ J$}
\loigiai{}
\end{ex}

\begin{ex}
Phát biểu nào sau đây là \textbf{sai} về sự phụ thuộc của nội năng của vật?
\choice
{phụ thuộc vào nhiệt độ vì nhiệt độ ảnh hưởng tới động năng}
{phụ thuộc vào cả nhiệt độ và thể tích}
{\True phụ thuộc vào nhiệt độ vì nhiệt độ ảnh hưởng tới thế năng}
{phụ thuộc vào thể tích vì thể tích ảnh hưởng tới thế năng}
\loigiai{}
\end{ex}

\begin{ex}
Đơn vị của nhiệt hoá hơi riêng là
\choice
{$kg/J$}
{\True $J/kg$}
{$J.kg$}
{$J$}
\loigiai{}
\end{ex}

\begin{ex}
Quá trình chuyển thể nào sau đây là quá trình tỏa nhiệt?
\choice
{\True Đông đặc, ngưng tụ}
{Hóa hơi, bay hơi}
{Nóng chảy, thăng hoa}
{Ngưng kết, thăng hoa}
\loigiai{}
\end{ex}

\begin{ex}
Phát biểu nào sau đây là \textbf{sai} khi nói về chất khí?
\choice
{Chất khí không có hình dạng và thể tích riêng}
{\True Các phân tử khí ở rất gần nhau}
{Lực tương tác giữa các phân tử khí rất yếu}
{Chất khí luôn luôn chiếm toàn bộ thể tích bình chứa và có thể nén được dễ dàng}
\loigiai{}
\end{ex}

\begin{ex}
Tại sao trên núi cao, ta \textbf{không} thể luộc chín trứng bằng nồi thông thường?
\choice
{\True nhiệt độ sôi giảm}
{nhiệt độ sôi tăng}
{nước không thể sôi}
{trên cao rất lạnh}
\loigiai{}
\end{ex}

\begin{ex}
Sự sôi là sự hóa hơi xảy ra ở
\choice
{\True cả trên bề mặt và trong lòng chất lỏng}
{trên bề mặt chất lỏng}
{trong lòng chất lỏng}
{phần chất lỏng tiếp xúc với bình chứa}
\loigiai{}
\end{ex}

\begin{ex}
Trong quá trình chất khí nhận nhiệt và sinh công thì $A$ và $Q$ trong biểu thức $\Delta U=Q+A$ có giá trị nào sau đây?
\choice
{$Q>0, A>0$}
{$Q<0, A>0$}
{\True $Q>0, A<0$}
{$Q<0, A<0$}
\loigiai{}
\end{ex}

\begin{ex}
Đặc điểm và tính chất nào dưới đây \textbf{không} liên quan đến chất rắn kết tinh?
\choice
{Có nhiệt độ nóng chảy xác định}
{Có cấu trúc tinh thể}
{Có dạng hình học xác định}
{\True Không có nhiệt độ nóng chảy xác định}
\loigiai{}
\end{ex}

\begin{ex}
Phát biểu nào sau đây là \textbf{sai}?
\choice
{\True Nội năng chính là nhiệt lượng của vật}
{Nội năng của một vật phụ thuộc vào nhiệt độ và thể tích của vật}
{Nội năng của vật có thể tăng hoặc giảm}
{Nội năng là một dạng năng lượng nên có thể chuyển hóa thành các dạng năng lượng khác}
\loigiai{}
\end{ex}

\begin{ex}
Sự truyền nhiệt là
\choice
{sự chuyển hóa năng lượng từ dạng này sang dạng khác}
{\True sự truyền trực tiếp nội năng từ vật này sang vật khác}
{sự chuyển hóa năng lượng từ nội năng sang dạng khác}
{sự truyền trực tiếp nội năng và chuyển hóa năng lượng từ dạng này sang dạng khác}
\loigiai{}
\end{ex}

\begin{ex}
Phát biểu nào sau đây \textbf{không} đúng? Tốc độ bay hơi của một lượng chất lỏng
\choice
{phụ thuộc vào áp suất của khí (hay hơi) trên bề mặt chất lỏng}
{\True không phụ thuộc vào bản chất của chất lỏng}
{càng lớn nếu diện tích bề mặt chất lỏng càng lớn}
{càng lớn nếu nhiệt độ chất lỏng càng cao}
\loigiai{}
\end{ex}

\begin{ex}
Khái niệm “ẩn nhiệt nóng chảy” được hiểu là phần năng lượng nhận thêm để
\choice
{giữ ổn định cấu trúc mạng tinh thể}
{làm tăng năng lượng liên kết giữa các phân tử}
{\True phá vỡ liên kết giữa các phân tử}
{làm tăng nhiệt độ của chất rắn khi nóng chảy}
\loigiai{}
\end{ex}

\begin{ex}
Khi nhiệt độ tuyệt đối tăng thêm $6\ K$ thì
\choice
{nhiệt độ Celsius tăng thêm $6^\circ C$}
{nhiệt độ Celsius giảm đi $267^\circ C$}
{\True nhiệt độ Celsius giảm đi $6^\circ C$}
{nhiệt độ Celsius tăng thêm $279^\circ C$}
\loigiai{}
\end{ex}

\begin{ex}
Cho biết nước đá có nhiệt nóng chảy riêng là $L=3,4.10^5\ J/kg$ và nhiệt dung riêng $c=2,09.10^3\ J/kg.K$. Nhiệt lượng cần cung cấp để làm nóng chảy cục nước đá khối lượng $50\ g$ và đang có nhiệt độ $-20^{\circ} C$ có giá trị bằng
\choice
{$104,5\ J$}
{$17\ kJ$}
{$2090\ J$}
{\True $19090\ J$}
\loigiai{}
\end{ex}

\begin{ex}
Có bốn vật A, B, C, D với các nhiệt độ ban đầu lần lượt là $t_A,\ t_B,\ t_C,\ t_D$. Cho vật A tiếp xúc với vật B và đến khi cân bằng nhiệt, nhiệt độ A giảm. Sau đó, vật A tiếp xúc với vật C đến khi cân bằng nhiệt, nhiệt độ A tăng. Tiếp theo, vật A tiếp xúc với vật D đến khi đạt cân bằng nhiệt thì nhiệt độ của A không đổi. Phát biểu nào sau đây là đúng?
\choice
{\True $t_B<t_A, t_C>t_A$, và $t_D=t_A$}
{$t_B>t_A, t_C<t_A$, và $t_D=t_A$}
{$t_B>t_A, t_C>t_A$, và $t_D<t_A$}
{$t_B<t_A, t_C<t_A$, và $t_D>t_A$}
\loigiai{}
\end{ex}

\begin{ex}
Đầu thép của một búa máy có khối lượng $12\ kg$ nóng lên thêm $20^\circ C$ sau $1,5$ phút hoạt động. Biết rằng chỉ có $30\%$ cơ năng của búa máy chuyển thành nhiệt năng ban đầu của đầu búa. Lấy nhiệt dung riêng của thép là $460\ J/kg.K$. Công suất trung bình của búa là
\choice
{$3670\ W$}
{$2245\ W$}
{\True $4089\ W$}
{$3067\ W$}
\loigiai{}
\end{ex}

\begin{ex}
Một thang đo nhiệt độ Z có nhiệt độ đóng băng của nước là $-25^\circ Z$ và khoảng cách mỗi độ chia trong thang đo độ K bằng $0,52$ lần khoảng cách mỗi độ chia trong thang đo nhiệt độ Z. Nhiệt độ sôi của nước trong thang nhiệt độ Z là
\choice
{$100^\circ Z$}
{$48^\circ Z$}
{$75^\circ Z$}
{\True $27^\circ Z$}
\loigiai{}
\end{ex}

\subsubsection*{PHẦN 2. Thí sinh trả lời từ câu 1 đến câu 4. Trong mỗi ý a), b), c), d) ở mỗi câu, thí sinh chọn đúng hoặc sai.}
\setcounter{ex}{0}
\begin{ex}%[2L1Y1-1]
Nghiên cứu mô hình cấu tạo chất.
\choiceTFt
{Ở thể khí, các phân tử ở rất gần nhau và chuyển động hỗn loạn không ngừng}
{Ở thể rắn, các phân tử dao động quanh vị trí cân bằng luôn thay đổi}
{\True Các phân tử chuyển động càng nhanh thì nhiệt độ của vật càng cao}
{Chất rắn kết tinh không có hình dạng và cấu trúc tinh thể xác định}
\loigiai{}
\end{ex}

%C:\Users\surface laptop\Desktop\TRONDE.tex
\begin{ex}%[2L1Y1-1]
Ở áp suất tiêu chuẩn, cho một ít nước đá ở dưới $0^\circ C$ vào bình chứa.
\choiceTFt
{Nhiệt độ nóng chảy của nước đá là $100^\circ C$}
{Trong suốt quá trình chuyển thành nước, nhiệt độ nước đá luôn tăng}
{\True Khi nước đá chuyển hoàn toàn thành nước, nếu vẫn tiếp tục đun nóng, nhiệt
độ của nước sẽ tăng dần và khi đạt $100^\circ C$, nước sẽ sôi và chuyển dần sang thể hơi}
{\True Đun nóng bình chứa thì nhiệt độ nước đá tăng dần, khi đạt nhiệt độ $0^\circ C$ thì nước đá tan dần thành nước}
\loigiai{}
\end{ex}

%C:\Users\surface laptop\Desktop\TRONDE.tex
\begin{ex}%[2L1B1-1]
\immini[thm]{Hình bên mô tả quan hệ giữa hai thang đo nhiệt độ X và Y
\choiceTFt
{\True Nhiệt độ $32^\circ Y$ sẽ tương ứng với nhiệt độ $0^\circ X$}
{\True Mối quan hệ giữa hai thang đo nhiệt độ được cho bởi công thức $T_Y=1,8T_X+32$}
{\True Tại nhiệt độ $-40$ độ thì giá trị trên hai thang đo là bằng nhau}
{Độ biến thiên nhiệt độ là $100^\circ X$ trên thang đo nhiệt độ X sẽ tương ứng với độ biến thiên $212^\circ Y$ trên thang đo nhiệt độ Y}
}{
\begin{tikzpicture}[draw=Mapcolor, very thick, scale = 0.8, line join=round, line cap=round, >=stealth]
\draw[->,line width = 1.2] (-1,0)--(4,0) node[above] {$^\circ X$};
\draw[->, line width = 1.2] (0,-1)--(0,5) node[right] {$^\circ Y$};
\draw[line width = 1.2] (-1,-1) -- (2.2,4.2);
\draw[fill=black] (0,0.6) circle (2pt) node [above left] {32};
\draw[dotted,line width = 1.2] (0,3.9) node [left] {212} -- (2,3.9) -- (2,0) node [below] {100} ;
\draw[fill=black]  (2,3.9) circle (2pt);
\draw (0,0) node [below left] {O};
\end{tikzpicture}}
\loigiai{}
\end{ex}

%C:\Users\surface laptop\Desktop\TRONDE.tex
\begin{ex}%[2L1Y1-1]
Nhiệt lượng là phần năng lượng truyền từ vật này sang vật khác do chênh lệch nhiệt độ giữa chúng.
\choiceTFt
{Một vật có nhiệt độ càng cao thì chứa càng nhiều nhiệt lượng}
{Nhiệt lượng là dạng năng lượng có đơn vị là jun}
{\True Trong sự truyền nhiệt không có sự chuyển hóa năng lượng từ dạng này sang dạng khác}
{\True Trong quá trình truyền nhiệt và thực hiện công, nội năng của vật không bảo toàn}
\loigiai{}
\end{ex}
\subsubsection*{PHẦN 3. Thí sinh trả lời từ câu 1 đến câu 6.}
\setcounter{ex}{0}
\begin{ex}%[2L1B1-1]
Bỏ một cục nước đá đang tan vào một nhiệt lượng kế chứa $1,5\ kg$ nước ở $30^\circ C$. Sau khi nó có cân bằng nhiệt người ta mang ra cân lại, khối lượng của nó chỉ còn lại $0,55\ kg$. Bỏ qua sự mất mát nhiệt. Biết $c_{\text{nước}}=4200\ J/kg.\text{độ};\ \lambda_{\text{nước đá}}=3,4.10^5\ J/kg$. Khối lượng của cục nước đá ban đầu là bao nhiêu kg (làm tròn kết quả đến hàng phần mười)?
\SA[oly]{1,1}
\loigiai{}
\end{ex}

\begin{ex}%[2L1K1-1]
\immini[thm]{Một bình nhiệt lượng kế ban đầu chứa $m_0=100\ g$ nước ở nhiệt độ $t_0=20^\circ C$. Người ta nhỏ đều đặn các giọt nước nóng vào nước đựng trong bình nhiệt lượng kế. Đồ thị biểu diễn sự phụ thuộc của nhiệt độ nước trong bình nhiệt lượng kế vào số giọt nước nóng nhỏ vào bình được biểu diễn ở đồ thị bên. Nhiệt độ của nước nóng là bao nhiêu độ Celsius?}
{\begin{tikzpicture}[draw=Mapcolor, very thick, scale = 0.8, line join=round, line cap=round, >=stealth]
\draw[->,line width = 1.2] (0,0)--(6.5,0) node[above] {N (giọt)};
\draw[->, line width = 1.2] (0,0)--(0,5) node[right] {t ($^o$C)};
\draw[fill = black] (0,0) circle(2pt) node [left] {\Large 0};
\draw[line width = 1.2] (0,1.8)..controls+(60:1)and+(-178:0.5)..(5,4);
\draw (0,2) node [left] {20};
\draw[line width = 1,dotted] (0,3) node [left] {30} -- (2,3) -- (2,0) node [below] {200};
\draw[line width = 1,dotted] (0,4) node [left] {40}-- (5,4) -- (5,0) node [below] {500};
\end{tikzpicture}}
\shortans[oly]{80}
\loigiai{}
\end{ex}

\begin{ex}%[2L1Y1-1]
Búa máy $10\ \text{tấn}$ rơi từ độ cao $2,3\ m$ xuống một cọc sắt ($c=460\ J/kg.K,\ m=200\ kg$). Biết $40\%$ động năng của búa biến thành nhiệt làm nóng cọc sắt. Cho $g=10\ m/s^2$. Búa máy rơi bao nhiêu lần thì cọc tăng nhiệt độ thêm $20^\circ C$?
\shortans[oly]{20}
\loigiai{}
\end{ex}

\begin{ex}%[2L1Y1-1]
Nhiệt lượng kế bằng đồng ($c_1=0,09\ cal/g.\text{độ}$) chứa nước ($c_2=1\ cal/g.\text{độ}$) ở $25^\circ C$. Khối lượng tổng cộng của nhiệt lượng kế là $475\ g$. Bỏ vào nhiệt lượng kế một vật bằng thau ($c_3=0,08\ cal/g.\text{độ}$) có khối lượng $400\ g$ ở $90^\circ C$. Nhiệt độ sau cùng của hệ khi cân bằng là $30^\circ C$. Khối lượng của nước là bao nhiêu gam?
\shortans[oly]{375}
\loigiai{}
\end{ex}

\begin{ex}%[2L1B1-1]
Có hai bình cách nhiệt. Bình I chứa $5\ \ell$ nước ở $60^\circ C$, bình II chứa $1\ \ell$ nước ở $20^\circ C$. Đầu tiên, rót một phần nước ở bình I sang bình II. Sau khi bình II cân bằng nhiệt, người ta lại rót từ bình II sang bình I một lượng nước bằng với lần rót trước. Nhiệt độ sau cùng của nước trong bình I là $59^\circ C$. Lượng nước đã rót từ bình này sang bình kia là bao nhiêu lít (làm tròn kết quả đến hàng phần trăm)?
\shortans[oly]{0,14}
\loigiai{}
\end{ex}

\begin{ex}%[2L1K1-1]
Trong hai bình cách nhiệt có chứa hai chất lỏng khác nhau ở hai nhiệt độ ban đầu khác nhau. Người ta dùng một nhiệt kế lần lượt nhúng đi nhúng lại vào bình 1, rồi vào bình 2. Chỉ số của nhiệt kế lần lượt là $40^\circ C,\ 8^\circ C,\ 39^\circ C,\ 9,5^\circ C$. Đến lần nhúng tiếp theo nhiệt kế chỉ bao nhiêu độ Celsius (làm tròn kết quả đến hàng phần mười)?
\shortans[oly]{38,1}
\loigiai{}
\end{ex}


